\section{Introduction}\label{sec:intro}

Viability theory asks whether a controller can keep a system inside a constraint set for as long as required \cite{Aubin1991}. Under full observation the computational side of that question is mature: viability kernels are approximated by backward recursions on grids \cite{SaintPierre1994}, by level-set formulations \cite{MitchellBayenTomlin2005}, and by Lagrangian reachability computations that scale considerably further \cite{Maidens2013}; set invariance has a well-developed convex theory \cite{Blanchini1999}; and the minimax reachability of target tubes under set-membership uncertainty was settled early \cite{BertsekasRhodes1971}. Under \emph{incomplete} observation the sufficiency side is carried by output-feedback conditions and estimation tubes \cite{Veliov1993}, but the necessity side---certificates whose validity \emph{proves} that no admissible controller can keep an information state safe---has lacked a computational bridge. The obstruction calculus of \cite{Abaee2026calc} characterises such certificates in continuous time; what is missing is a finite object, computable by a standard solver and checkable by a third party, that carries the same force.

The difficulty is one of quantifier order. The controller announces a policy adapted to its observation history; the adversary then chooses the hidden mode, the initial state inside its cell, and a measurable disturbance. A certificate of nonviability must therefore cover \emph{every} measurable, information-adapted policy. Discretising the control signal and showing that the resulting finite program is infeasible proves only that no \emph{held} control is safe, which is a strictly weaker statement; the direction of approximation matters, and getting it wrong makes infeasibility-by-discretisation arguments unsound.

\subsection{Contributions}

This paper builds the bridge and the tooling around it, and it does so under an explicit division of labour: the mathematics proves soundness and completeness, the solver computes, and exact rational witnesses---not the solver---certify what is printed.

\begin{enumerate}[leftmargin=1.6em,itemsep=2pt]
\item \textbf{A continuous-to-finite theorem for information states} (Theorem~\ref{thm:bridge}). The admissible controls are approximated from \emph{outside} by their block moments, a relaxation that restricts nothing, and the adversary from \emph{inside} by finitely many stored labels and scenarios. The resulting finite linear program has an optimal value $\rhoHG$ that is a certified lower bound on the continuous-time safety value $\J$, and a certified inflation that is an upper bound:
\[
\rhoHG\;\le\;\J\;\le\;\rhoHG+\max_a\ebar_a+\delta_\beta+Lh_t .
\]
Every dual-feasible point of the program is a finite obstruction certificate valid against every measurable policy; nonviability is detected along every admissible refinement sequence; and the value has an exact atomic dual representation, proved through Sion's minimax theorem on the mixed extension of the label set (Corollary~\ref{cor:dual}).
\item \textbf{Certificate size is a matter of information--time rank, not input dimension} (Section~\ref{sec:size}). Whenever a positive certificate exists, one exists with at most $r+1$ safety labels, $r$ the rank of the finite safety-row matrix; a construction with a \emph{scalar} input produces obstructions whose minimal certificates need arbitrarily many witnesses (Proposition~\ref{prop:rank}); and the \emph{obstruction depth} of a belief---the smallest number of compatible modes whose joint presence is already nonviable---is a matching lower bound (Proposition~\ref{prop:depth}). The Helly-type intuition ``$m+1$ compatible states suffice'' is false for common blind control.
\item \textbf{An instance carried exactly from threshold to solver} (Section~\ref{sec:instance}). On a three-branch delayed-observation instance every pair of compatible states is viable, with explicit pairwise policies, and a common blind control keeps every branch safe throughout the blind window; nevertheless the triple is obstructed past the exact delay $\tau^\ast=7/50$. The finite value obeys the closed law $\rho(\calH)=\tfr{3}{50}-\tfr14\sum_{J\in\calH}|J|^2$ on \emph{every} mesh aligned with the observation time (Proposition~\ref{prop:meshlaw}), certified by one mesh-independent primal--dual witness pair; a dual-simplex solver reproduces it to $4.2\times10^{-17}$ across a refinement campaign, and the smallest certifying program has 39 variables.
\item \textbf{The certificate is prescriptive} (Section~\ref{sec:ladder}). Its weights and facet multipliers decode into the instantaneous shared braking authorities of each sub-belief, an exact delay hierarchy $7/50<(15-\sqrt{183})/6<(10-\sqrt{58})/6$ (Theorem~\ref{thm:pairs}), and an \emph{erosion identity} that prices candidate observers: the same witness loses a computable margin when an information cell is split, the ranking of observers by erosion coincides with the ranking of pairs by delay tolerance, and for an observer installed at an intermediate time the identity brackets the critical installation time within $[3/25,(6-\sqrt{29})/5]$, an interval of width below $3\times10^{-3}$ (Proposition~\ref{prop:observer}).
\item \textbf{Exact verification by construction} (Section~\ref{sec:computation}). Primal and dual witnesses are checked in rational arithmetic; a residual-bounded test makes the certificate robust to enclosed data; the Farkas core is mechanised; and a verification script re-derives all 63 exact statements of the paper.
\end{enumerate}

What is \emph{not} claimed is equally explicit (Section~\ref{sec:related}): the two-sided sandwich idea as such belongs to the reachability literature \cite{Maidens2013}, moment relaxation as such to \cite{Lasserre2001,Parrilo2003,Lasserre2008}, and scenario discretisation as such to \cite{CalafioreCampi2006,CampiGaratti2008}. The novelty lies in assembling these into a certificate that is valid against every measurable information-adapted policy under partial observation, in the rank and depth results, and in the prescriptive reading.

\subsection{Organisation}

Section~\ref{sec:setting} fixes the model, the information structure and the safety value. Section~\ref{sec:certificates} defines finite certificates and states the bridge theorem. Section~\ref{sec:size} treats certificate size. Section~\ref{sec:instance} carries the instance, the mesh law and the solver campaign. Section~\ref{sec:ladder} reads the certificate prescriptively. Section~\ref{sec:extensions} records sound extensions beyond exogenous mode observation. Section~\ref{sec:computation} covers complexity and verification, Section~\ref{sec:related} the related work, and Section~\ref{sec:discussion} concludes. Appendix~\ref{app:bridge} proves the bridge theorem; Appendix~\ref{app:ladder} proves the ladder results. The Supplementary Material (SM) contains the remaining proofs, the derivation of the moduli from primitive data, the complete instance arithmetic, the observer-pricing details, and the verification record.
